\section{Kinematics, symmetry and conventions}\label{sec:tt-kinematics}
Let $k$ denote the relative momentum in the $t+t$ center-of-mass frame and $\mu$ the reduced mass. The Sommerfeld parameter for repulsive Coulomb interactions is
\begin{equation}\label{eq:sommerfeld}
\eta=\frac{\alpha Z_1Z_2\mu}{k}, \qquad Z_1=Z_2=1,
\end{equation}
in units $\hbar=c=1$. Every tabulated value in a later analysis must state whether energy means laboratory kinetic energy, center-of-mass kinetic energy, or excitation energy relative to a specified threshold.

\subsection{Identical-fermion selection rules}
The total state must be antisymmetric under interchange of the two tritons. In the elastic two-cluster spin basis, singlet spin $S=0$ is antisymmetric and therefore permits even orbital angular momenta $\ell$; triplet spin $S=1$ is symmetric and therefore permits odd $\ell$. In particular, the elastic $S$ wave is spin singlet. These selection rules are essential before comparing phase shifts or low-energy parameters between publications.

\subsection{Coulomb-modified effective-range expansion}
For an isolated elastic $S$ wave, one standard convention introduces the nuclear phase shift measured relative to Coulomb functions, $\delta_0^C$, through
\begin{equation}\label{eq:coulomb-ere}
C_\eta^2\,k\cot\delta_0^C(k)+2k\eta\,h(\eta)
=-\frac{1}{a_C}+\frac{r_C}{2}k^2+\mathcal{O}(k^4),
\end{equation}
where $C_\eta^2=2\pi\eta/(\ee^{2\pi\eta}-1)$ and $h(\eta)=\RePart\psi(\ii\eta)-\ln\eta$ (equivalently the real part of the standard Coulomb function). Here $a_C$ and $r_C$ are \emph{convention-dependent Coulomb-modified} parameters, not automatically the scattering length and range of a hypothetical short-range-only interaction. Equation~\eqref{eq:coulomb-ere} does not justify applying a single-channel fit across an inelastic threshold.

% CHECK: In a full revision, state explicitly the convention for the Coulomb phase sigma_l and nuclear amplitude normalization.
